\documentclass[11pt]{article}

\usepackage[margin=1in]{geometry}
\usepackage{amsmath,amssymb,amsthm}
\usepackage{graphicx}
\usepackage[hidelinks]{hyperref}

\newtheorem{theorem}{Theorem}
\newtheorem{lemma}{Lemma}
\newtheorem{remark}{Remark}

\newcommand{\R}{\mathbb R}
\newcommand{\conv}{\operatorname{conv}}
\newcommand{\lin}{\operatorname{lin}}
\newcommand{\polar}{^\circ}

\title{A Planar Equality Case Outside the Source-Style Pattern\\
for the Dual Log-Brunn-Minkowski Inequality}
\author{\texttt{agent\_lane\_03}}
\date{2026-07-18}

\begin{document}
\maketitle

\section*{Source and Target}

The source paper is Christos Saroglou,
\emph{More on logarithmic sums of convex bodies}, arXiv:1409.4346
\cite{Saroglou2016}.  In Section 6, Saroglou proves the dual logarithmic
Brunn-Minkowski inequality
\[
 \left|(\lambda\cdot K+_0(1-\lambda)\cdot L)\polar\right|
 \leq |K\polar|^\lambda |L\polar|^{1-\lambda},
 \qquad 0<\lambda<1.
\]
Immediately after the theorem, the paper says that it is plausible that the
equality cases are the same as those in Conjecture 1.1, with non-symmetric
bodies allowed, but does not address this.

\begin{figure}[h]
\centering
\includegraphics[width=0.96\linewidth]{figures/open_problem_crop.png}
\caption{Theorem 6.1 and the equality-case comment in arXiv:1409.4346, p. 17.}
\end{figure}

The construction below gives equality in the dual inequality for a planar pair
which is neither homothetic nor of the natural product/dilate form suggested
by the equality cases in Conjecture 1.1.  Thus the source-style equality-case
analogy is false in this literal sense.

\section*{Definitions}

For a convex body \(C\subset\R^2\) containing the origin in its interior, write
\[
   C\polar=\{y\in\R^2: x\cdot y\leq 1\ \text{for all }x\in C\}
\]
for its polar.  The logarithmic sum is the Wulff body determined by the
support-function inequality
\[
 \lambda\cdot K+_0(1-\lambda)\cdot L
 =
 \{y\in\R^2: y\cdot u\leq h_K(u)^\lambda h_L(u)^{1-\lambda}
      \text{ for all } u\in\R^2\}.
\]

\section*{Proof Intuition}

Saroglou's proof of the dual inequality reduces the volume estimate to
log-convexity of the volume of polytopes whose vertices move by independent
exponential radial factors.  In the plane, the area of a polygon is the sum
of determinants of adjacent vertices.  If the radial exponents alternate
\(0,1,0,1\) around a quadrilateral, every adjacent determinant has exponent
\(1\).  The area is therefore exactly log-affine.

Taking polars turns this log-affine moving quadrilateral into a pair
\(K,L\).  The logarithmic sum has polar containing the intermediate moving
quadrilateral, while the dual log-Brunn-Minkowski inequality gives the
opposite volume bound.  Hence equality is forced for the true logarithmic sum,
not merely for a discrete approximation.  The chosen quadrilateral is
deliberately not a product polar, so this equality falls outside the natural
source-style product/dilate list.

\section*{The Counterexample}

\begin{theorem}
Let
\[
 x_1=(2,0),\qquad x_2=(0,1),\qquad
 x_3=(-1,\tfrac12),\qquad x_4=(0,-1),
\]
and, for \(0\leq t\leq 1\), set
\[
   P_t=\conv\{x_1,e^t x_2,x_3,e^t x_4\}.
\]
Put \(K=P_0\polar\) and \(L=P_1\polar\).  Then, for every
\(0<\lambda<1\),
\[
 \left|(\lambda\cdot K+_0(1-\lambda)\cdot L)\polar\right|
 =
 |K\polar|^\lambda |L\polar|^{1-\lambda}.
\]
Nevertheless, \(K\) is not a two-dimensional Cartesian product of one
dimensional factors, even after a linear change of coordinates, and \(L\) is
not a dilate of \(K\).  Hence this equality pair is outside the natural
planar product/dilate equality pattern suggested by the source comment.
\end{theorem}

\begin{proof}
First note that the four listed points occur in counterclockwise order and
that \(P_t\) contains the origin in its interior for \(0\leq t\leq1\).  Indeed,
with \(q=e^t\geq1\), direct edge-orientation checks give
\[
2q,\quad q,\quad q,\quad 2q
\]
for the signed areas of the triangles formed by the origin and consecutive
vertices, and the edge-cross-products of the polygon are
\[
3q-1,\quad 2q,\quad 3q+1,\quad 4q,
\]
all positive.

The area of \(P_t\) is especially simple.  By the shoelace formula,
\[
\begin{aligned}
2|P_t|
&=\det(x_1,qx_2)+\det(qx_2,x_3)
  +\det(x_3,qx_4)+\det(qx_4,x_1)  \\
&=2q+q+q+2q=6q.
\end{aligned}
\]
Thus
\[
   |P_t|=3e^t.                                      \tag{1}
\]

Let \(a_1=a_3=0\) and \(a_2=a_4=1\), so that
\[
   P_t=\conv\{e^{a_i t}x_i:1\leq i\leq4\}.
\]
Since each \(x_i\) is an exposed vertex of \(P_0=K\polar\), its polar
constraint is active and
\[
   h_K(x_i)=1.
\]
Likewise \(e^{a_i}x_i\) is the corresponding exposed vertex of
\(P_1=L\polar\), so
\[
   h_L(x_i)=e^{-a_i}.
\]

Fix \(0<\lambda<1\) and write
\[
   M_\lambda=\lambda\cdot K+_0(1-\lambda)\cdot L.
\]
The defining support inequalities for the logarithmic sum, tested only in the
four directions \(x_i\), imply
\[
   M_\lambda
   \subseteq
   Q_\lambda
   :=
   \{y\in\R^2: y\cdot x_i\leq e^{-(1-\lambda)a_i},
       \ 1\leq i\leq4\}.
\]
Taking polars reverses inclusion, and the polar of \(Q_\lambda\) is
\[
   Q_\lambda\polar
   =
   \conv\{e^{(1-\lambda)a_i}x_i:1\leq i\leq4\}
   =
   P_{1-\lambda}.
\]
Hence
\[
   P_{1-\lambda}\subseteq M_\lambda\polar.           \tag{2}
\]

Saroglou's dual logarithmic Brunn-Minkowski inequality gives
\[
   |M_\lambda\polar|
   \leq |K\polar|^\lambda |L\polar|^{1-\lambda}
   = |P_0|^\lambda |P_1|^{1-\lambda}.
\]
Using (1),
\[
   |P_0|^\lambda |P_1|^{1-\lambda}
   =3^\lambda (3e)^{1-\lambda}
   =3e^{1-\lambda}
   =|P_{1-\lambda}|.
\]
Together with (2), this forces equality:
\[
   |M_\lambda\polar|=|P_{1-\lambda}|
   =|K\polar|^\lambda |L\polar|^{1-\lambda}.
\]

It remains to verify that this is not a source-style planar equality case.
In dimension two, a Cartesian product of one-dimensional factors, after a
linear change of coordinates, has a polar whose vertices split into two pairs
lying on two lines through the origin.  The vertices of \(P_0=K\polar\) do not
admit such a split.  The only origin-collinear pair among the four vertices is
\(\{x_2,x_4\}\), while
\[
   \det(x_1,x_3)=1\neq0.
\]
Thus \(K\) is not a two-dimensional product body of one-dimensional factors.

Finally, \(L\) is not a dilate of \(K\).  If \(L=cK\), then
\[
   P_1=L\polar=c^{-1}K\polar=c^{-1}P_0.
\]
But \(P_1\) has the two \(y\)-axis vertices \(e x_2,e x_4\), whereas \(P_0\)
has \(x_2,x_4\), so any scalar relation would force \(c^{-1}=e\).  The same
scalar would then have to send \(x_1\) to \(e x_1\), but \(x_1\) itself is a
vertex of \(P_1\).  This contradiction proves that \(L\) is not a dilate of
\(K\).
\end{proof}

\section*{Verification Notes}

The script \texttt{code/check\_quadrilateral\_equality.py} verifies the exact
determinants
\[
   2,\ 1,\ 1,\ 2
\]
which give \(|P_t|=3e^t\), checks that \(P_0\) has only one origin-collinear
pair of vertices, and records the elementary obstruction to a scalar homothety
\(P_1=cP_0\).

Command used:
\begin{quote}
\texttt{conda run --no-capture-output -n sandbox python}\\
\texttt{code/check\_quadrilateral\_equality.py}
\end{quote}

This computation is not a substitute for the proof; it is a reproducible
sanity check for the coordinates used in the construction.

\section*{Scope and Novelty Check}

This packet does not refute Saroglou's dual log-Brunn-Minkowski inequality.
It refutes the equality-case analogy if that analogy is read as requiring the
same planar product/dilate pattern as in the log-Brunn-Minkowski equality
statement.  It also does not classify all equality cases in the dual theorem.

The run's cheap indexes were searched for arXiv:1409.4346, the source title,
``log-Brunn'', ``dual log-Brunn-Minkowski'', ``equality cases'',
``B-conjecture'', ``cone-volume'', and related log-Minkowski terms.  No prior
packet for this paper or route was found.

A bounded web search on 2026-07-18 for exact and nearby phrases including
``dual log-Brunn-Minkowski equality cases'', ``Saroglou dual B-conjecture
equality cases polytopes'', and the source title found the source article and
later general log-Brunn-Minkowski literature, but no existing equality-case
classification or counterexample for this Section 6 comment.

\section*{Human Review Recommendation}

Recommended verdict: likely valid counterexample to the equality-case analogy,
with narrow scope.  Reviewers should mainly check the intended interpretation
of the source's informal equality-case comment: the mathematical equality
above is exact, while the classification conclusion being refuted is the
natural source-style product/dilate pattern rather than a formal theorem.

\begin{thebibliography}{9}

\bibitem{Saroglou2016}
Christos Saroglou,
\emph{More on logarithmic sums of convex bodies},
arXiv:1409.4346; \emph{Mathematika} 62 (2016), 818--841.

\end{thebibliography}

\end{document}
